% 4.2 第七章 电流与磁场.tex
\chapter{第七章 电流与磁场}

\section{恒定电流和恒定电场 电动势}
    \subsection{电流形成的条件}
    定义单位时间内通过导线任一截面的电荷量为电流。

    在导体中形成电流的条件是：（1）导体（1）
    \subsection{恒定电流和恒定电场}
    \subsection{电流和电流密度}
    定义单位时间内通过垂直于电流方向单位面积的电量为导体中某点电流密度矢量的$\boldsymbol{j}$大小，$\boldsymbol{j}$的方向与正电荷在该点
    \subsection{欧姆定律的微分形式 焦耳—楞次定律的微分形式}
    \subsection{电源及电源电动势}
    电源电动势在量值

\section{恒定磁场和磁感应强度}
    \subsection{磁性起源于电荷的运动}
    \subsection{磁场 磁感应强度\texorpdfstring{$\boldsymbol{B}$}{B}}

\section{毕奥‑萨伐尔定律}
    \subsection{毕奥‑萨伐尔定律}
    在真空中
    \subsection{毕奥‑萨伐尔定律应用举例}
    \subsection{匀速运动电荷的磁场}

\section{磁场中的高斯定理}
    \subsection{磁感应线}
    \subsection{磁通量}
    \subsection{恒定磁场中的高斯定理}

\section{真空中恒定磁场的安培环路定理}
    \subsection{恒定磁场的安培环路定理}
    \subsection{安培环路定理的应用}

\section{磁场对运动电荷和载流导线的作用}
    \subsection{洛伦兹力}
    \subsection{带电粒子在磁场中的运动}
    \subsection{应用电场和磁场控制带电粒子运动的实例}
    \subsection{安培力}

\section{磁力的功}
    \subsection{磁力对运动载流导线的功}
    \subsection{磁力矩对转动载流线圈的功}
    磁力矩对转动载流线圈的功等于电流强度与回路包围面积内的磁通量的增量的乘积。

\section{磁介质中的恒定磁场}
    \subsection{磁介质及其磁化}
    \subsection{磁介质中的高斯定理和安培环路定理}
    \subsection{铁磁质}

\endinput